% Generated draft: requires mathematical and source review.
\documentclass[11pt]{article}
\usepackage[margin=25mm]{geometry}
\usepackage{fontspec}
\setmainfont{Latin Modern Roman}
\usepackage{amsmath,amssymb,mathtools,unicode-math}
\setmathfont{Latin Modern Math}
\usepackage{enumitem,needspace,xurl,hyperref}
\setlength{\parindent}{0pt}
\setlength{\parskip}{5pt}
\setlength{\emergencystretch}{3em}
\newcommand{\N}{\mathbb{N}}
\newcommand{\Z}{\mathbb{Z}}
\newcommand{\Q}{\mathbb{Q}}
\newcommand{\R}{\mathbb{R}}
\newcommand{\C}{\mathbb{C}}
\newcommand{\F}{\mathbb{F}}
\newcommand{\A}{\mathbb{A}}
\newcommand{\PP}{\mathbb{P}}
\newcommand{\E}{\mathbb{E}}
\newcommand{\eps}{\varepsilon}
\newcommand{\dd}{\,\mathrm{d}}
\begin{document}
\section*{UnsolvedMath ID: 20002868}
\textbf{Source problem code:} catalogue entry
\subsection{Definitions and mathematical statement}
Let $n\geq1$ and let $R(z,\bar z)\in\C[z_1,\ldots,z_n,\bar z_1,\ldots,\bar z_n]$ be real-valued for all $z\in\C^n$, equivalently Hermitian under conjugating coefficients and interchanging $z,\bar z$. D'Angelo's problem asks for necessary and sufficient conditions for
\[
\exists N\in\Z_{\geq1}\ \exists k\in\Z_{\geq1}\ \exists p_1,\ldots,p_k\in\C[z_1,\ldots,z_n],\qquad
R(z,\bar z)^N=\sum_{j=1}^k|p_j(z)|^2
\]
as a polynomial identity. The exponent is positive; allowing $N=0$ would trivialize the question. The zero polynomial is permitted in this general question. Negative real constants and negative multiples of squared norms can qualify using even powers, so the problem must not silently impose $R\geq0$ or an odd exponent.

The supplied partial results are as follows. If $R$ has holomorphic and antiholomorphic degrees at most $d$, set
\[
\widetilde R(z_0,z;\bar z_0,\bar z)
=z_0^d\bar z_0^dR(z/z_0,\bar z/\bar z_0),
\]
interpreted as its polynomial expansion, bihomogeneous of bidegree $(d,d)$. For each fixed $N\geq1$, $R^N$ is a squared norm of holomorphic polynomials if and only if $\widetilde R^N$ is such a squared norm; the homogeneous summands on the latter side have degree $dN$. The leading bidegree corner $R_{d,d}$ also has a squared-norm $N$th power whenever $R$ does, with the zero corner allowed in this statement.

Polarize $R$ to the polynomial $\mathcal R(z,\zeta)$ in independent complex variables. If $R(a,\bar a)=0$ and $R$ has a squared-norm power, the aid asserts
\[
\mathcal R\in\mathfrak m_a(z)\,\mathfrak m_{\bar a}(\zeta),
\]
where $\mathfrak m_a(z)=(z_1-a_1,\ldots,z_n-a_n)$ and the second ideal is defined analogously. In particular $\mathcal R(a,\zeta)\equiv0$ and $\mathcal R(z,\bar a)\equiv0$.

For the one-variable consequence assume explicitly $R\not\equiv0$ and bidegree exactly $(d,d)$. Each distinct diagonal zero $a\in\C$, meaning $R(a,\bar a)=0$, forces the factor $(z-a)(\zeta-\bar a)$ in the polarization. Consequently there are at most $d$ distinct diagonal zeros. Within this class, equality holds precisely when
\[
R(z,\bar z)=c\prod_{\nu=1}^d|z-a_\nu|^2,
\qquad c\in\R\setminus\{0\},\quad a_1,\ldots,a_d\text{ pairwise distinct}.
\]
For $d=0$ the product is empty and the statement concerns a nonzero real constant. Distinctness is essential to the equality count: a general degree-$d$ polynomial with repeated roots does not give $d$ distinct zeros. The supplied zero-count criterion and homogenization reductions do not give necessary and sufficient conditions for all Hermitian polynomials, and they do not provide an exponent bound.
\subsection{Short English statement}
The question asks which real-valued complex polynomials have some positive power equal to a finite sum of squared moduli of holomorphic polynomials. The supplied progress preserves this property under bihomogenization and forces strong factors at diagonal zeros. In one variable a nonzero bidegree-$(d,d)$ example has at most $d$ distinct zeros, with equality only in the stated distinct-root squared-modulus form. A full classification is not supplied.
\subsection{Sources}
\begin{itemize}
\item Local registry: permanent catalogue entry, exact source code catalogue entry; \texttt{aim-several-complex-variables-notes.json}, \texttt{notes[25]}.
\item AIM workshop list \emph{Emerging applications of complexity for CR mappings}, Problem 2 attributed to J. D'Angelo, \url{https://aimath.org/WWN/crmappings/crmappings.pdf}. The original remarks and literature field are empty; no complete authoritative publication citation is supplied.
\item The verification field says the extraction lost the exponent in $R^N$ and placement of the summation indices, but ends before reproducing the repaired original formula. The displayed positive-power interpretation follows the supplied research summary.
\item Supplied research attempt 1, registry updated 14 August 2026: partial result, open-at-run label, low novelty confidence, machine-generated and not peer reviewed. It supplies homogenization, leading-corner and zero-factor obstructions, and the one-variable equality case without an external citation.
\end{itemize}
\end{document}
